\subsection{Associative Algebra Structures}

% Role: foundational algebra operations, regular representations, and
% endomorphism-algebra examples.

The module viewpoint in Definition \ref{def:operator-module-and-matrix-module} encodes one operator at a time.\\
An algebra records a vector space of compatible products or operators.\\
The guiding example is $\LL(\mathcal{V})$, whose vector-space operations are pointwise and whose multiplication is composition.

\begin{definition} \label{def:algebra-over-field} \hlt{Algebra over a Field}\\
Let $\mathbb{F}$ be a field. An \hlt{$\mathbb{F}$-algebra} is an $\mathbb{F}$-vector space $\mathcal{A}$ with multiplication $\mathcal{A}\times\mathcal{A}\to\mathcal{A}$, $(a,b)\mapsto ab$.\\
The structure $(\mathcal{A},+,\cdot)$ is a ring. Scalar compatibility is $(\forall\,\lambda\in\mathbb{F})(\forall\,a,b\in\mathcal{A})\,\lambda(ab)=(\lambda a)b=a(\lambda b)$.
\end{definition}

\begin{definition} \label{def:unital-associative-algebra} \hlt{Unital Algebras and Basic Properties}\\
Let $\mathcal{A}$ be an $\mathbb{F}$-algebra.
\begin{enumerate}[label=(\roman*)]
\setlength{\itemsep}{0pt}
\item $\mathcal{A}$ is a \hlt{unital associative $\mathbb{F}$-algebra} if its underlying ring is equipped with an identity $1_{\mathcal{A}}$.
\item $\mathcal{A}$ is \hlt{finite-dimensional} iff $\dim_{\mathbb{F}}\mathcal{A}<\infty$.
\item $\mathcal{A}$ is \hlt{commutative} iff $(\forall\,a,b\in\mathcal{A})\,ab=ba$.
\item If $\mathcal{A}$ is unital, then $a\in\mathcal{A}$ is \hlt{invertible} iff $(\exists\,b\in\mathcal{A})\,ab=ba=1_{\mathcal{A}}$.
\end{enumerate}
\end{definition}

\begin{convention} \label{conv:algebra-convention-relative-to-ring-definition} \hlt{Algebra Convention}\\
For a fixed field $\mathbb{F}$, \hlt{algebra} means unital associative $\mathbb{F}$-algebra unless stated otherwise.
\end{convention}

\begin{example} \label{ex:basic-associative-algebras} \hlt{Basic Algebras}
\begin{enumerate}[label=(\roman*)]
\setlength{\itemsep}{0pt}
\item If $\mathbb{F}\subseteq\mathbb{E}$ are fields, then $\mathbb{E}$ is an $\mathbb{F}$-algebra by $r\cdot x=rx$ and multiplication in $\mathbb{E}$.
\item $\mathbb{F}[x]$ is a commutative $\mathbb{F}$-algebra with unit $1$ and scalar embedding $r\mapsto r$.
\item If $n\geq1$, then $\operatorname{Mat}(n,n,\mathbb{F})$ is an $\mathbb{F}$-algebra with unit $I_n$ and product $AB$.
\item If $\mathcal{V}$ is an $\mathbb{F}$-vector space, then $\LL(\mathcal{V})=\operatorname{End}_{\mathbb{F}}(\mathcal{V})$ is an $\mathbb{F}$-algebra. Its unit is $I_{\mathcal{V}}$, and its product is composition $S\circ T$.
\end{enumerate}
\end{example}

\begin{definition} \label{def:subalgebra-generated-subalgebra} \hlt{Subalgebras and Generated Subalgebras}\\
Let $\mathcal{A}$ be an $\mathbb{F}$-algebra.
\begin{enumerate}[label=(\roman*)]
\setlength{\itemsep}{0pt}
\item A \hlt{subalgebra}, written $\mathcal{B}\leq_{\mathbb{F}\text{-alg}}\mathcal{A}$, is an $\mathbb{F}$-linear subspace $\mathcal{B}\leq_{\mathbb{F}}\mathcal{A}$ containing $1_{\mathcal{A}}$. It satisfies $(\forall\,b,b'\in\mathcal{B})\,bb'\in\mathcal{B}$.
\item For $X\subseteq\mathcal{A}$, the \hlt{subalgebra generated by $X$} is $\langle X\rangle_{\operatorname{subalg}}\defeq\bigcap\limits\{\mathcal{B}\leq_{\mathbb{F}\text{-alg}}\mathcal{A}:X\subseteq\mathcal{B}\}$.
\end{enumerate}
\end{definition}

\begin{definition} \label{def:center-and-central-algebra} \hlt{Center and Central Algebras}
\begin{enumerate}[label=(\roman*)]
\setlength{\itemsep}{0pt}
\item The \hlt{center} of an $\mathbb{F}$-algebra $\mathcal{A}$ is $Z(\mathcal{A})\defeq\{a\in\mathcal{A}:(\forall\,x\in\mathcal{A})\,ax=xa\}$.
\item $\mathcal{A}$ is \hlt{central over $\mathbb{F}$} iff $Z(\mathcal{A})=\mathbb{F}1_{\mathcal{A}}\defeq\{r1_{\mathcal{A}}:r\in\mathbb{F}\}$.
\end{enumerate}
\end{definition}

\begin{proposition} \label{prop:field-embeds-in-center-of-algebra} \hlt{Centers and Central Scalar Embeddings}\\
Let $\mathcal{A}$ be an $\mathbb{F}$-algebra.
\begin{enumerate}[label=(\roman*)]
\setlength{\itemsep}{0pt}
\item $Z(\mathcal{A})\leq_{\mathbb{F}\text{-alg}}\mathcal{A}$.
\item If $1_{\mathcal{A}}\neq0$, then $\iota_{\mathcal{A}}:\mathbb{F}\hookrightarrow Z(\mathcal{A})$, $r\mapsto r1_{\mathcal{A}}$, is an injective unital ring homomorphism.
\item If $R$ is a unital ring and $\mathbb{F}\leq Z(R)$ is a subfield, then $R$ is an $\mathbb{F}$-algebra via $r\cdot x\defeq rx$.
\end{enumerate}
\begin{proof}
\begin{enumerate}[label=(\roman*)]
\setlength{\itemsep}{0pt}
\item Identity belongs to $Z(\mathcal{A})$. If $(r,s,a,b,x)\in\mathbb{F}^2\times Z(\mathcal{A})^2\times\mathcal{A}$, then $(ra+sb)x=r(ax)+s(bx)=r(xa)+s(xb)=x(ra+sb)$.\\
Thus $Z(\mathcal{A})\leq_{\mathbb{F}}\mathcal{A}$. For $a,b\in Z(\mathcal{A})$ and $x\in\mathcal{A}$, $(ab)x=a(bx)=a(xb)=(ax)b=(xa)b=x(ab)$. Thus $Z(\mathcal{A})\leq_{\mathbb{F}\text{-alg}}\mathcal{A}$.
\item Set $\iota\defeq\iota_{\mathcal{A}}$. For $r,s\in\mathbb{F}$, $\iota(r+s)=\iota(r)+\iota(s)$, $\iota(rs)=\iota(r)\iota(s)$, and $\iota(1)=1_{\mathcal{A}}$ by the algebra axioms.\\
For $(r,a)\in\mathbb{F}\times\mathcal{A}$, $\iota(r)a=(r1_{\mathcal{A}})a=ra=a(r1_{\mathcal{A}})=a\iota(r)$, so $\operatorname{im}\iota\subseteq Z(\mathcal{A})$.\\
If $\iota(r)=0$ and $r\neq0$, then $1_{\mathcal{A}}=\iota(r^{-1})\iota(r)=0$, contradiction; hence $\ker\iota=\{0\}$.
\item Set $r\cdot x\defeq rx$. Then $R$ is an $\mathbb{F}$-vector space by the field laws in $\mathbb{F}\subseteq R$.\\
Scalar compatibility $(\forall\,r\in\mathbb{F})(\forall\,x,y\in R)\,r\cdot(xy)=(r\cdot x)y=x(r\cdot y)$ holds by $\mathbb{F}\subseteq Z(R)$. Thus $R$ is an $\mathbb{F}$-algebra.
\end{enumerate}
\hlt{Structural Commentary}: Center closure supports the scalar embedding, while a central subfield reconstructs the scalar action.\\
\hlt{Key Insight}: Scalar multiplication is multiplication by the corresponding central element.
\end{proof}
\end{proposition}

\begin{definition} \label{def:algebra-ideal-simple-quotient} \hlt{Ideals, Simple Algebras, and Quotients}\\
Let $\mathcal{A}$ be an $\mathbb{F}$-algebra.
\begin{enumerate}[label=(\roman*)]
\setlength{\itemsep}{0pt}
\item A \hlt{two-sided ideal}, written $\mathcal{I}\triangleleft\mathcal{A}$, is an $\mathbb{F}$-linear subspace $\mathcal{I}\leq_{\mathbb{F}}\mathcal{A}$. It satisfies $(\forall\,a\in\mathcal{A})(\forall\,u\in\mathcal{I})\,au,ua\in\mathcal{I}$.
\item For $X\subseteq\mathcal{A}$, the \hlt{ideal generated by $X$} is $\langle X\rangle_{\operatorname{ideal}}\defeq\bigcap\limits\{\mathcal{I}\triangleleft\mathcal{A}:X\subseteq\mathcal{I}\}$.
\item $\mathcal{A}$ is \hlt{simple} iff $\mathcal{A}\neq\{0_{\mathcal{A}}\}$ and its only two-sided ideals are $\{0_{\mathcal{A}}\}$ and $\mathcal{A}$.
\item If $\mathcal{I}\triangleleft\mathcal{A}$, \hlt{quotient algebra} $\mathcal{A}/\mathcal{I}$ is quotient vector space. Product and identity are $(a+\mathcal{I})(b+\mathcal{I})\defeq ab+\mathcal{I}$ and $1_{\mathcal{A}/\mathcal{I}}\defeq1_{\mathcal{A}}+\mathcal{I}$.
\end{enumerate}
\end{definition}

\begin{proposition} \label{prop:generated-subalgebra-and-ideal-formulas} \hlt{Concrete Generated Substructures}\\
Let $\mathcal{A}$ be an $\mathbb{F}$-algebra and let $X\subseteq\mathcal{A}$. Set $\mathcal{W}_X\defeq\{1_{\mathcal{A}}\}\cup\{x_1\cdots x_m:m\in\N,\ (x_1,\ldots,x_m)\in X^m\}$.\\
Set $\mathcal{I}_X\defeq\{\sum\limits_{j=1}^{n}a_jx_jb_j:n\in\N_0,\ (a_j,b_j,x_j)_{j=1}^{n}\in(\mathcal{A}\times\mathcal{A}\times X)^n\}$, with the empty sum at $n=0$.
\begin{enumerate}[label=(\roman*)]
\setlength{\itemsep}{0pt}
\item $\langle X\rangle_{\operatorname{subalg}}=\operatorname{span}_{\mathbb{F}}(\mathcal{W}_X)$.
\item $\langle X\rangle_{\operatorname{ideal}}=\mathcal{I}_X$.
\end{enumerate}
\begin{proof}
\begin{enumerate}[label=(\roman*)]
\setlength{\itemsep}{0pt}
\item Products of members of $\mathcal{W}_X$ belong to $\mathcal{W}_X$, so $\operatorname{span}_{\mathbb{F}}(\mathcal{W}_X)$ is a subalgebra containing $X$ by bilinearity.\\
If $\mathcal{B}\leq_{\mathbb{F}\text{-alg}}\mathcal{A}$ and $X\subseteq\mathcal{B}$, then $\mathcal{W}_X\subseteq\mathcal{B}$ and $\operatorname{span}_{\mathbb{F}}(\mathcal{W}_X)\subseteq\mathcal{B}$. The equality follows by Definition \ref{def:subalgebra-generated-subalgebra}(ii).
\item $\mathcal{I}_X\leq_{\mathbb{F}}\mathcal{A}$ by closure of finite sums under $\mathbb{F}$-linear combinations, and $X\subseteq\mathcal{I}_X$ by $x=1_{\mathcal{A}}x1_{\mathcal{A}}$.\\
For $(a,b)\in\mathcal{A}^2$ and $u=\sum\limits_{j=1}^{n}a_jx_jb_j\in\mathcal{I}_X$, both $au=\sum\limits_{j=1}^{n}(aa_j)x_jb_j$ and $ub=\sum\limits_{j=1}^{n}a_jx_j(b_jb)$ belong to $\mathcal{I}_X$.\\
Thus $\mathcal{I}_X\triangleleft\mathcal{A}$. If $\mathcal{J}\triangleleft\mathcal{A}$ and $X\subseteq\mathcal{J}$, then $(\forall\,a,b\in\mathcal{A})(\forall\,x\in X)\,axb\in\mathcal{J}$, so $\mathcal{I}_X\subseteq\mathcal{J}$. The equality follows by Definition \ref{def:algebra-ideal-simple-quotient}(ii).
\end{enumerate}
\hlt{Structural Commentary}: Closure of words constructs the generated subalgebra; two-sided products construct the generated ideal.\\
\hlt{Key Insight}: Intersections agree with the explicit closures forced by their defining operations.
\end{proof}
\end{proposition}

\begin{definition} \label{def:algebra-homomorphism} \hlt{Algebra Homomorphisms}\\
Let $\mathcal{A},\mathcal{B}$ be $\mathbb{F}$-algebras.
\begin{enumerate}[label=(\roman*)]
\setlength{\itemsep}{0pt}
\item An \hlt{algebra homomorphism} is an $\mathbb{F}$-linear map $\sigma\in\LL(\mathcal{A},\mathcal{B})$. It satisfies $(\forall\,a,b\in\mathcal{A})\,\sigma(ab)=\sigma(a)\sigma(b)$ and $\sigma(1_{\mathcal{A}})=1_{\mathcal{B}}$.
\item An algebra homomorphism is a \hlt{monomorphism} iff it is injective.
\item An algebra homomorphism is an \hlt{isomorphism} iff it is bijective.
\end{enumerate}
\end{definition}

\begin{proposition} \label{prop:algebra-kernel-first-isomorphism} \hlt{Kernel, Image, and First Isomorphism Theorem}\\
Let $\sigma:\mathcal{A}\to\mathcal{B}$ be an $\mathbb{F}$-algebra homomorphism and set $K\defeq\ker\sigma$.
\begin{enumerate}[label=(\roman*)]
\setlength{\itemsep}{0pt}
\item $K\triangleleft\mathcal{A}$.
\item $\operatorname{im}\sigma\leq_{\mathbb{F}\text{-alg}}\mathcal{B}$.
\item The map $\overline{\sigma}:\mathcal{A}/K\to\operatorname{im}\sigma$, $a+K\mapsto\sigma(a)$, is an $\mathbb{F}$-algebra isomorphism.
\end{enumerate}
\begin{proof}
\begin{enumerate}[label=(\roman*)]
\setlength{\itemsep}{0pt}
\item $K\leq_{\mathbb{F}}\mathcal{A}$ by linearity. If $(a,u)\in\mathcal{A}\times K$, then $\sigma(au)=\sigma(a)\sigma(u)=0=\sigma(u)\sigma(a)=\sigma(ua)$, so $au,ua\in K$ and $K\triangleleft\mathcal{A}$.
\item The image is an $\mathbb{F}$-linear subspace by linearity and contains $1_{\mathcal{B}}=\sigma(1_{\mathcal{A}})$.\\
It is closed under multiplication by $(\forall\,a,b\in\mathcal{A})\,\sigma(a)\sigma(b)=\sigma(ab)$. Thus $\operatorname{im}\sigma\leq_{\mathbb{F}\text{-alg}}\mathcal{B}$.
\item Define $\overline{\sigma}(a+K)\defeq\sigma(a)$. The equivalence $a+K=a'+K\Longleftrightarrow a-a'\in K$ holds by Proposition \ref{prop:quotient-module-well-defined}(i).\\
The equivalence $a-a'\in K\Longleftrightarrow\sigma(a)=\sigma(a')$ holds by $K=\ker\sigma$.\\
Thus $\overline{\sigma}$ is well defined and injective; it is surjective by its codomain.\\
For $(a,a',r)\in\mathcal{A}^2\times\mathbb{F}$, $\overline{\sigma}((a+K)+(a'+K))=\sigma(a)+\sigma(a')$ and $\overline{\sigma}(r(a+K))=r\sigma(a)$.\\
Multiplicativity is $\overline{\sigma}((a+K)(a'+K))=\sigma(a)\sigma(a')$.\\
The identity $\overline{\sigma}(1_{\mathcal{A}}+K)=1_{\mathcal{B}}$ completes the verification that $\overline{\sigma}$ is an $\mathbb{F}$-algebra isomorphism.
\end{enumerate}
\hlt{Structural Commentary}: Kernel absorption forms the quotient, while image closure supplies the induced map's algebra codomain.\\
\hlt{Key Insight}: Cosets agree exactly when their representatives have the same image.
\end{proof}
\end{proposition}

\begin{definition} \label{def:left-regular-representation} \hlt{Left Regular Representation}\\
Let $\mathcal{A}$ be an $\mathbb{F}$-algebra. For $a\in\mathcal{A}$, the \hlt{left multiplication map} is $L_a:\mathcal{A}\to\mathcal{A}$, $x\mapsto ax$.\\
The map $L_{\mathcal{A}}:\mathcal{A}\to\LL(\mathcal{A})$, $a\mapsto L_a$, is the \hlt{left regular representation} of $\mathcal{A}$.
\end{definition}

\begin{theorem} \label{thm:left-regular-representation-algebra} \hlt{Left Regular Representation}\\
If $\mathcal{A}$ is an $\mathbb{F}$-algebra, then $L_{\mathcal{A}}:\mathcal{A}\hookrightarrow\LL(\mathcal{A})$, $a\mapsto L_a$, is an $\mathbb{F}$-algebra monomorphism.
\begin{proof}
For $(a,x,y,r)\in\mathcal{A}^3\times\mathbb{F}$, $L_a(x+y)=L_a(x)+L_a(y)$ and $L_a(rx)=rL_a(x)$, so $L_a\in\LL(\mathcal{A})$.\\
For $(a,b,r,x)\in\mathcal{A}^3\times\mathbb{F}$, $L_{a+b}(x)=L_a(x)+L_b(x)$ and $L_{ra}(x)=rL_a(x)$. Hence $L_{\mathcal{A}}$ is $\mathbb{F}$-linear by Assumption \ref{assump:function-extensionality}.\\
For $a,b,x\in\mathcal{A}$, $L_{ab}(x)=a(bx)=(L_aL_b)(x)$ and $L_{1_{\mathcal{A}}}(x)=x$. Hence $L_{ab}=L_aL_b$ and $L_{1_{\mathcal{A}}}=I_{\mathcal{A}}$ by Assumption \ref{assump:function-extensionality}.\\
If $L_a=0$, then $a=L_a(1_{\mathcal{A}})=0$, so $L_{\mathcal{A}}$ is injective.\\
\hlt{Structural Commentary}: Bilinearity defines the algebra map, while unit evaluation proves injectivity.\\
\hlt{Key Insight}: The identity recovers each element by $L_a(1_{\mathcal{A}})=a$.
\end{proof}
\end{theorem}

Figure \ref{fig:algebra-left-regular-representation} records the multiplicativity and unit-evaluation identities used in the proof.

\begin{illustrationfigure}[H]
  \centering
  \begin{tikzpicture}[
  >=Stealth,
  x=0.9cm,
  every node/.style={font=\small},
  alg/.style={draw=blue!60!black, fill=blue!7, rounded corners=2pt,
    minimum width=2.2cm, minimum height=0.64cm, align=center, inner sep=4pt},
  op/.style={draw=dartmouthgreen!65!black, fill=dartmouthgreen!8, rounded corners=2pt,
    minimum width=2.65cm, minimum height=0.64cm, align=center, inner sep=4pt},
  vec/.style={draw=gray!55, fill=gray!6, rounded corners=2pt,
    minimum width=1.8cm, minimum height=0.56cm, align=center, inner sep=3pt},
  arr/.style={->, line width=0.75pt, draw=black!70},
  note/.style={font=\scriptsize, inner sep=2pt, fill=white}
]
  \node[alg] (A) at (0,0) {$a\in\mathcal{A}$};
  \node[op] (L) at (4.15,0) {$L_a\in\LL(\mathcal{A})$};
  \node[vec] (x) at (8.05,0) {$x\in\mathcal{A}$};
  \node[vec] (ax) at (11.45,0) {$ax\in\mathcal{A}$};

  \draw[arr] (A) -- node[note, above] {$a\mapsto L_a$} (L);
  \draw[arr] (L) -- node[note, above] {acts on} (x);
  \draw[arr] (x) -- node[note, above] {$L_a$} (ax);
\end{tikzpicture}

  \caption{The left regular representation sends $a\in\mathcal{A}$ to $L_a:x\mapsto ax$; the identities $L_{ab}=L_aL_b$ and $L_a(1_{\mathcal{A}})=a$ encode multiplicativity and injectivity.}
  \label{fig:algebra-left-regular-representation}
\end{illustrationfigure}

\begin{computationmethod} \label{meth:structure-constants-algebra} \hlt{Building Multiplication from Structure Constants}\\
\hlt{Input}: A vector space $V=\bigoplus\limits_{i\in I}\mathbb{F}b_i$ and coefficients $c_{ij}^k\in\mathbb{F}$. The support condition is $(\forall\,i,j\in I)\,\abs{\{k\in I:c_{ij}^k\neq0\}}<\infty$.\\
\hlt{Output}: A bilinear multiplication on $V$, together with coordinate tests for associativity, commutativity, and a basis identity.
\begin{enumerate}[label=(\roman*)]
\setlength{\itemsep}{0pt}
\item Define $b_ib_j\defeq\sum\limits_{k\in I}c_{ij}^kb_k$. Extend bilinearly by $(\sum\limits_i r_ib_i)(\sum\limits_j s_jb_j)\defeq\sum\limits_{i,j,k}r_is_jc_{ij}^kb_k$.
\item The product is associative iff $(\forall\,i,j,k,m\in I)\,\sum\limits_{\ell\in I}c_{ij}^{\ell}c_{\ell k}^{m}=\sum\limits_{\ell\in I}c_{jk}^{\ell}c_{i\ell}^{m}$.
\item The product is commutative iff $(\forall\,i,j,k\in I)\,c_{ij}^k=c_{ji}^k$.
\item A basis vector $b_e$ is an identity iff $(\forall\,j,m\in I)\,c_{ej}^m=\delta_j^m=c_{je}^m$.
\end{enumerate}
\end{computationmethod}

\begin{definition} \label{def:group-algebra} \hlt{Group Algebra}\\
For a group $G$ and field $\mathbb{F}$, the \hlt{group algebra} is $\mathbb{F}[G]\defeq\{\sum\limits_{g\in G}r_gg:r_g\in\mathbb{F},\ \abs{\{g\in G:r_g\neq0\}}<\infty\}$.\\
Its vector-space operations are coefficientwise, and $(\sum\limits_g r_gg)(\sum\limits_h s_hh)\defeq\sum\limits_{g,h}r_gs_h(gh)$.
\end{definition}

\begin{proposition} \label{prop:group-algebra-structure} \hlt{Structure of a Group Algebra}\\
Let $G$ be a group and let $\mathbb{F}$ be a field.
\begin{enumerate}[label=(\roman*)]
\setlength{\itemsep}{0pt}
\item $G$ is an $\mathbb{F}$-basis of $\mathbb{F}[G]$.
\item With the multiplication in Definition \ref{def:group-algebra}, $\mathbb{F}[G]$ is an $\mathbb{F}$-algebra with identity $e_G$.
\item $\mathbb{F}[G]$ is commutative iff $G$ is abelian.
\item If $G$ is finite, then $\dim_{\mathbb{F}}\mathbb{F}[G]=\abs{G}$.
\end{enumerate}
\begin{proof}
\begin{enumerate}[label=(\roman*)]
\setlength{\itemsep}{0pt}
\item[$(i)$] $G$ is an $\mathbb{F}$-basis by the unique finite-support coefficient expansions in Definition \ref{def:group-algebra}.
\item[$(ii)$] The product is bilinear by construction. The two associative products are trilinear and agree on basis triples by $(gh)k=g(hk)$.\\
The identity relation follows by bilinearity and $(\forall\,g\in G)\,e_Gg=g=ge_G$.
\item[$(iii)$] ($\Rightarrow$) If $\mathbb{F}[G]$ is commutative, then $(\forall\,g,h\in G)\,gh=hg$, so $G$ is abelian.
\item[$(iii)$] ($\Leftarrow$) If $G$ is abelian, then $\mathbb{F}[G]$ is commutative by bilinearity and $(\forall\,g,h\in G)\,gh=hg$.
\item[$(iv)$] The dimension equals the cardinality of the basis $G$ by part (i).
\end{enumerate}
\hlt{Structural Commentary}: Basis multiplication transfers the group law to the algebra, while basis equality detects commutativity.\\
\hlt{Key Insight}: The group elements remain linearly independent inside $\mathbb{F}[G]$.
\end{proof}
\end{proposition}

\begin{example} \label{ex:cyclic-group-algebra-regular-matrices} \hlt{Cyclic Group Algebras and Circulant Shifts}\\
Let $G=\langle\alpha\rangle$ have order $n$, and let $\mathcal{B}=(1,\alpha,\ldots,\alpha^{n-1})$ be the ordered $\mathbb{F}$-basis of $\mathbb{F}[G]$.\\
With exponents read modulo $n$, $(\forall\,j,k\in\{0,\ldots,n-1\})\,L_{\alpha^k}(\alpha^j)=\alpha^{j+k}$, so $L_{\alpha^k}$ is the $k$-fold cyclic shift of $\mathcal{B}$.\\
Set $P\defeq[L_{\alpha}]_{\mathcal{B}}$. Then $(\forall\,k\in\{0,\ldots,n-1\})\,[L_{\alpha^k}]_{\mathcal{B}}=P^k$. For $n\geq5$, the following coordinate form displays the cyclic shift.
\begin{align*}
P&=
\begin{bmatrix}
0&0&\cdots&0&1\\
1&0&\cdots&0&0\\
0&1&\ddots&0&0\\
\vdots&\vdots&\ddots&\vdots&\vdots\\
0&0&\cdots&1&0
\end{bmatrix}.
\end{align*}
Thus $[L_{\sum\limits_{k=0}^{n-1}c_k\alpha^k}]_{\mathcal{B}}=\sum\limits_{k=0}^{n-1}c_kP^k$, and these are precisely the circulant matrices.
\end{example}

\begin{proposition} \label{prop:endomorphism-algebra-matrix-algebra-isomorphism} \hlt{Endomorphisms and Matrices Are the Same Algebra after Choosing a Basis}\\
Let $\dim\mathcal{V}=n\geq1$ and fix an ordered basis $\mathcal{B}$ of $\mathcal{V}$. Then $\mu_{\mathcal{B}}:\LL(\mathcal{V})\to\operatorname{Mat}(n,n,\mathbb{F})$, $T\mapsto[T]_{\mathcal{B}}$, is an $\mathbb{F}$-algebra isomorphism.
\begin{proof}
By Theorem \ref{thm:linear-maps-and-matrices-same-data}, $\mu_{\mathcal{B}}$ is an $\mathbb{F}$-linear bijection. For $S,T\in\LL(\mathcal{V})$, $\mu_{\mathcal{B}}(ST)=\mu_{\mathcal{B}}(S)\mu_{\mathcal{B}}(T)$ by Proposition \ref{prop:composition-and-matrix-multiplication}.\\
Since $\mu_{\mathcal{B}}(I_{\mathcal{V}})=I_n$, $\mu_{\mathcal{B}}$ is an $\mathbb{F}$-algebra isomorphism.\\
\hlt{Structural Commentary}: Linear coordinate encoding becomes algebra isomorphism by composition compatibility, identity preservation.\\
\hlt{Key Insight}: A fixed basis converts composition exactly into matrix multiplication.
\end{proof}
\end{proposition}

\begin{theorem} \label{thm:endomorphism-algebra-center-ideals-simple} \hlt{Center and Ideals of an Endomorphism Algebra}\\
Assume Zorn's lemma, and let $0\neq\mathcal{V}$ be an $\mathbb{F}$-vector space. Set $\mathcal{F}(\mathcal{V})\defeq\{T\in\LL(\mathcal{V}):\dim\operatorname{im}T<\infty\}$.
\begin{enumerate}[label=(\roman*)]
\setlength{\itemsep}{0pt}
\item $Z(\LL(\mathcal{V}))=\mathbb{F}I_{\mathcal{V}}$, so $\LL(\mathcal{V})$ is central.
\item $\mathcal{F}(\mathcal{V})\triangleleft\LL(\mathcal{V})$.
\item $(\forall\,\mathcal{J}\triangleleft\LL(\mathcal{V}))\,[\mathcal{J}\neq\{0\}\Rightarrow\mathcal{F}(\mathcal{V})\subseteq\mathcal{J}]$.
\item $\LL(\mathcal{V})$ is simple iff $\dim\mathcal{V}<\infty$.
\end{enumerate}
\begin{proof}
\begin{enumerate}[label=(\roman*)]
\setlength{\itemsep}{0pt}
\item[$(i)$] For $(v,\varphi)\in\mathcal{V}\times\mathcal{V}'$, write $v\otimes\varphi\in\LL(\mathcal{V})$ for $x\mapsto\varphi(x)v$. Fix $T\in Z(\LL(\mathcal{V}))$.\\
For $0\neq v\in\mathcal{V}$, choose $\varphi\in\mathcal{V}'$ with $\varphi(v)=1$ by Proposition \ref{prop:functionals-separate-points}(iii), with $\{0\}\leq\mathcal{V}$.\\
Set $\lambda_v\defeq\varphi(Tv)$. Then $Tv=T(v\otimes\varphi)(v)=(v\otimes\varphi)T(v)=\lambda_vv$. Fix $0\neq v_0\in\mathcal{V}$. If $v=cv_0\neq0$, then $\lambda_v=\lambda_{v_0}$ by linearity.\\
If $v,v_0$ are independent, then $\lambda_{v+v_0}(v+v_0)=\lambda_vv+\lambda_{v_0}v_0$, so $\lambda_v=\lambda_{v_0}$ by independence. Hence $T=\lambda_{v_0}I_{\mathcal{V}}$ by Assumption \ref{assump:function-extensionality}.\\
Since scalar operators commute with every endomorphism, $Z(\LL(\mathcal{V}))=\mathbb{F}I_{\mathcal{V}}$.
\item[$(ii)$] The zero map belongs to $\mathcal{F}(\mathcal{V})$. For $(r,s,T,U)\in\mathbb{F}^2\times\mathcal{F}(\mathcal{V})^2$, $\operatorname{im}(rT+sU)\subseteq\operatorname{im}T+\operatorname{im}U$.\\
Thus $\mathcal{F}(\mathcal{V})\leq_{\mathbb{F}}\LL(\mathcal{V})$. For $R,S\in\LL(\mathcal{V})$ and $T\in\mathcal{F}(\mathcal{V})$, $\operatorname{im}(RTS)\subseteq R(\operatorname{im}T)$; $\mathcal{F}(\mathcal{V})\triangleleft\LL(\mathcal{V})$.
\item[$(iii)$] Let $\{0\}\neq\mathcal{J}\triangleleft\LL(\mathcal{V})$. Choose $0\neq S\in\mathcal{J}$, $e\in\mathcal{V}$ with $Se\neq0$. Choose $\eta\in\mathcal{V}'$ with $\eta(Se)=1$ by Proposition \ref{prop:functionals-separate-points}(iii), $\{0\}\leq\mathcal{V}$.\\
For $(v,\varphi)\in\mathcal{V}\times\mathcal{V}'$, $v\otimes\varphi=(v\otimes\eta)S(e\otimes\varphi)\in\mathcal{J}$, so $\mathcal{J}$ contains every rank-one map.\\
If $T\in\mathcal{F}(\mathcal{V})$ and $(w_1,\ldots,w_m)$ is a basis of $\operatorname{im}T$, let $\psi_j$ be its coordinate functionals.\\
For $x\in\mathcal{V}$, $T(x)=\sum\limits_{j=1}^{m}\psi_j(Tx)w_j$. The identity $T=\sum\limits_{j=1}^{m}w_j\otimes(\psi_j\circ T)$ holds by Assumption \ref{assump:function-extensionality}.\\
Each summand belongs to $\mathcal{J}$, so $T\in\mathcal{J}$ and $\mathcal{F}(\mathcal{V})\subseteq\mathcal{J}$.
\item[$(iv)$] ($\Rightarrow$) Suppose $\LL(\mathcal{V})$ is simple, $\dim\mathcal{V}=\infty$. Choose $0\neq v\in\mathcal{V}$, $\varphi\in\mathcal{V}'$ with $\varphi(v)=1$ by Proposition \ref{prop:functionals-separate-points}(iii), $\{0\}\leq\mathcal{V}$.\\
Then $0\neq v\otimes\varphi\in\mathcal{F}(\mathcal{V})$, whereas $I_{\mathcal{V}}\notin\mathcal{F}(\mathcal{V})$; $\mathcal{F}(\mathcal{V})$ is nonzero proper ideal by part (ii), contradiction. Hence $\dim\mathcal{V}<\infty$.
\item[$(iv)$] ($\Leftarrow$) If $\dim\mathcal{V}<\infty$, then $\mathcal{F}(\mathcal{V})=\LL(\mathcal{V})$, so every nonzero two-sided ideal equals $\LL(\mathcal{V})$ by part (iii).\\
Since $I_{\mathcal{V}}\neq0$, the algebra $\LL(\mathcal{V})$ is nonzero and therefore simple.
\end{enumerate}
\hlt{Structural Commentary}: Rank-one probes force scalar central maps, while rank-one factorizations control every nonzero ideal.\\
\hlt{Key Insight}: Sandwiching one nonzero operator produces every rank-one operator.
\end{proof}
\end{theorem}

\begin{perspective} \label{per:algebra-chapter-operator-connection} \hlt{Operator Connection}\\
$T\in\LL(\mathcal{V})$ determines an $\mathbb{F}[z]$-module by $z\cdot v=T(v)$.\\
The left regular map $L_{\mathcal{A}}:\mathcal{A}\hookrightarrow\LL(\mathcal{A})$ identifies $\mathcal{A}$ with a subalgebra of its endomorphism algebra and supplies a faithful representation.
\end{perspective}
